\section{Geometric limitations of compressing the existing witnesses}
\label{sec:geometric-compression-limits}

The bound of $41$ total coordinates is not a lower bound for
definitions of $\mathbb Z$.  This section proves a more limited
statement: certain proposed ways to compress the existing elliptic
point modules cannot work.  A new arithmetic construction can avoid
these limitations by restricting or replacing those modules while
preserving only the input projection.

\subsection{Finite algebra requires geometric finiteness}

\begin{proposition}\label{prop:finite-algebra-dimension-obstruction}
Let $V$ be a geometrically irreducible variety of dimension $d$ over
a field $F$, and let $S\subseteq V(F)$ be Zariski dense.  Suppose
there are finitely many varieties $Y_j$, morphisms
$\pi_j:Y_j\longrightarrow\mathbb A_F^m$ with finite geometric
fibres, and rational maps $g_j:Y_j\dashrightarrow V$, defined
where they are used, such that
\[
                         S\subseteq\bigcup_jg_j(Y_j(F)).
\]
Then $m\ge d$.
\end{proposition}
\begin{proof}
The dimension theorem gives $\dim Y_j\le m$ because each
geometric fibre of $\pi_j$ is finite.  Every irreducible component
of the closure of $g_j(Y_j)$ therefore has dimension at most $m$.
If $m<d$, the union of these finitely many closures is a proper
closed subset of the irreducible variety $V$.  It cannot contain
the Zariski dense set $S$.
\end{proof}

In particular, a finite triangular monic algebra over $m$ rational
parameters, followed by rational decoding formulas, cannot cover
a Zariski dense $d$-dimensional family when $m<d$.  A finite list
of exceptional charts does not change this conclusion.

This observation also distinguishes finite rational fibres from
finite geometric fibres.  Pasten proves that, for a positive-rank
elliptic curve over a number field, there is an affine open set
$U$ and a morphism $U\times U\to\mathbb A^1$ which is injective
on rational points, although those rational points are Zariski
dense in the surface \cite[Corollary 1.2]{P3Pasten2020}.
Its rational fibres have at most one point.  Its generic geometric
fibre has dimension one.  Proposition~\ref{prop:finite-algebra-dimension-obstruction}
shows that this injection does not supply a finite-algebra decoder
for the pair from one rational coordinate.

\subsection{The rank-one point group is dense in a surface}

We now establish the obstruction for the exact elliptic curve
used in the construction.  The proof needs no height estimate,
prime-pattern theorem, or restricted converse for elliptic curves.

\begin{lemma}\label{lem:conjugate-twist-density}
Let $F$ be a number field, let $K/F$ be quadratic with involution
$\sigma$, and let $d\in K^*$.  Put $N=d\sigma(d)\in F^*$ and
suppose $\sqrt N\notin K(i)$.  On
\[
                         E_d:y^2=x^3-d^2x
\]
let $P\in E_d(K)$ be nontorsion.  The cyclic group generated by
$P$, viewed in the Weil restriction
$A=\operatorname{Res}_{K/F}E_d$, is geometrically Zariski dense
in the two-dimensional variety $A$.
\end{lemma}
\begin{proof}
After extension to an algebraic closure,
$A\simeq E_d\times E_{\sigma(d)}$.  Put $u=d/\sqrt N$.
The change of variables
\[
 E_{\sigma(d)}\longrightarrow E_d,
 \qquad (x,y)\longmapsto(u^2x,u^3y)
\]
is an isomorphism because
$u^4\sigma(d)^2=d^2$.  Write $R$ for the image of
$\sigma(P)$.  There is an automorphism $\tau$ fixing $K(i)$
and changing the sign of $\sqrt N$.  Consequently
\[
                            \tau(P)=P,\qquad\tau(R)=-R.
\]

The geometric endomorphism ring of $E_d$ is $\mathbb Z[i]$,
and all its endomorphisms are defined over $K(i)$.  One can see
this from complex uniformization: the order-four automorphism
$(x,y)\mapsto(-x,iy)$ makes the period lattice a rank-one
$\mathbb Z[i]$-module; since $\mathbb Z[i]$ is a principal ideal
domain, its multiplier ring is exactly $\mathbb Z[i]$.
In particular $\tau$ fixes every such endomorphism.

Suppose $\alpha(P)+\beta(R)=0$ for endomorphisms
$\alpha,\beta\in\mathbb Z[i]$.  Applying $\tau$ gives
$\alpha(P)-\beta(R)=0$.  Thus $2\alpha(P)=2\beta(R)=0$.
A nonzero endomorphism of an elliptic curve is an isogeny and
has finite kernel.  Since $P$ and $R$ are nontorsion, it follows
that $\alpha=\beta=0$.

For completeness, this absence of relations implies the claimed
density.  The Zariski closure of a subgroup of an algebraic group
is an algebraic subgroup.  If the closure of $\mathbb Z(P,R)$
were proper in $E_d^2$, its identity component would have
dimension zero or one.  Dimension zero would make $(P,R)$
torsion.  In dimension one, let $C$ be that identity component.
Both projections $C\to E_d$ are nonzero isogenies: otherwise
a nonzero multiple of $P$ or $R$ would vanish.  If these maps
are $\phi,\psi$, a dual isogeny to $\phi$ gives a nonzero
relation
\[
                       [\deg\phi]y-\alpha x=0
                       \quad\hbox{on }C
\]
for an endomorphism $\alpha$ of $E_d$.  A nonzero multiple of
$(P,R)$ lies in $C$, so this is a nontrivial endomorphism
relation between $P$ and $R$, contradicting the previous
paragraph.  The standard isogeny facts used here may be found
in \cite[Sections 8 and 12]{P3MilneAbelian}.
\end{proof}

\begin{corollary}[The point modules used here]
\label{cor:point-modules-dimension}
For $K=\mathbb Q(\sqrt2)$, $d=3-\sqrt2$, and the explicit
point
\[
 P=(57/8-7\sqrt2/2,-165/8+417\sqrt2/32),
\]
the set $\mathbb ZP$ is Zariski dense in
$\operatorname{Res}_{K/\mathbb Q}E_d$.  Every nonzero
finite-index subgroup has the same property.  Nine independent
copies of the point module consequently form a Zariski dense
subset of a variety of dimension eighteen.

A construction which covers all those independent point choices
using finitely many rational decoding charts over a finite
algebra needs at least eighteen retained rational parameters.
In particular, replacing each point module separately by one
rational parameter and finitely many algebraic choices is
impossible.
\end{corollary}
\begin{proof}
Here $N(d)=7$ and the quadratic subfields of
$\mathbb Q(\sqrt2,i)$ are
$\mathbb Q(\sqrt2)$, $\mathbb Q(i)$, and
$\mathbb Q(\sqrt{-2})$.  Thus $\sqrt7\notin K(i)$.

The stated point is nontorsion by an elementary local calculation.
The curve has good reduction at three, the reduction of $P$ is
the nonzero two-torsion point $(\sqrt2,0)$, and
\[
 z(2P)=
 \frac{3396356159088-4135204317756\sqrt2}{3287811625151},
 \qquad z=-x/y,
\]
has valuation one at the unramified quadratic place above three.
For a point of positive integral depth there, multiplication by
a three-adic unit preserves depth, while multiplication by three
increases depth by exactly one.  Indeed the integral formal
group has odd multiplication series, whose first term is $nz$;
the higher terms in $[3]z$ have valuation at least $3v(z)$,
strictly greater than $1+v(z)$.  No nonzero integer can therefore
annihilate $2P$.  The displayed exact calculation is reproduced
by \texttt{check\_elliptic\_density\_data.py}.

Apply Lemma~\ref{lem:conjugate-twist-density}.  Replacing $P$
by a nonzero multiple preserves its nontorsion property and
the same argument.  A Cartesian product of Zariski dense sets
is dense in the product variety, which has dimension eighteen.
The final assertion follows from
Proposition~\ref{prop:finite-algebra-dimension-obstruction}.
\end{proof}

\paragraph{Scope of this obstruction.}
The nine point choices are not independent after every arithmetic
condition of the final definition has been imposed.  The corollary
therefore gives a limitation on separate replacement of the point
modules and on witness-preserving finite-algebra parametrization.
It is not a lower bound of eighteen witnesses for $\mathbb Z$,
and it does not exclude a new definition with at most ten total
coordinates.  Such a reduction must exploit additional arithmetic
relations among the points, replace the point construction, or use
a different kind of representation.
